%% kapa: exam=1 round=36 year=2025 period=1 form=A subject=economics subject_ko=경제학원론 seq=19 number=59
%% source: Q-Net 감정평가사 기출문제 (cst00309) artlSeq=5231525 file=raw/1008/5231525/2025년 제36회 감정평가사 1차시험 1교시 문제지 원본.pdf
%% conversion: kordoc 4.24.51 -> tools/gen1.py (auto); review=proofread-sonnet; flags=none
\begin{question}{59}
두 명의 수요자로 구성된 $X$ 재 시장에 관한 설명으로 옳은 것은? (단, $P$ 는 가격, $Q_{1}$ 과 $Q_{2}$ 는 각각 개인 1과 2의 수요함수)
\par\noindent\fbox{\parbox{\dimexpr\linewidth-2\fboxsep-2\fboxrule\relax}{\vspace{1pt}
$\circ$ $Q_{1}=10-\frac{1}{2}P$ , $\quad Q_{2}=20-P$ \par
$\circ$ 한계비용은 $16$ 으로 일정하다.
\vspace{1pt}}}
\begin{choices}
  \item $X$재가 사적재일 경우 경쟁시장에서의 균형거래량은 10이다.
  \item $X$재가 사적재일 경우 경쟁시장에서의 균형거래량은 8이다.
  \item $X$재가 공공재일 경우 최적 생산량은 8이다.
  \item $X$재가 공공재일 경우 최적 생산량은 10이다.
  \item 재화의 성격과 관계없이 균형거래량과 최적 생산량은 10이다.
\end{choices}
\end{question}
\answer{3}
